\scanpage{207}

or to the pairs $(\alpha,\beta)$ of elements of $\widehat{\Pi}_{0,3}$
satisfying the loop relation (37)
\[
 (97)=(37)\qquad
 \unsure{(\alpha\rho(\alpha)^{-1})}
 =(\beta\sigma_\infty(\beta)^{-1})\ell_1^{\unsure{\nu}}
\]
(which we will then complete by taking $f=\beta$).

\margin{$\alpha,\beta$ \unsure{determined} \unsure{modulo}
$(\alpha,\beta)\longmapsto
(\ell_0^{\unsure{n}}\alpha,\ell_0^{\unsure{n}}\beta)$.}

Up to changing $\alpha$ to $\ell_0^\nu\alpha=\underline{\alpha}_0$,
these were the conventions of \unsure{§ 44|§ 4 h} (where the formula
\unsure{thus} takes the form
\[
 \underbrace{\underline{\alpha}_0\rho(\underline{\alpha}_0)^{-1}}_{\xi}
 =\ell_0^\nu
 \underbrace{(f\sigma_\infty(f)^{-1})}_{\unsure{\varphi=\eta}}
\]
\unsure{).}

Taking (79) into account, the conditions

\deleted{\uline{2° Case}
\unsure{(98)} $g=1$ i.e. $\unsure{\alpha_0}=1$ i.e. $f=\alpha$,
i.e. $\xi=f\rho(f)^{-1}$ i.e.}

(93) (94) also take the form
\[
 (98)\qquad u_{\alpha,\beta,f}
 =\mu(t\sigma_\infty+u)\tau_\infty^{\alpha_2},
\]
and imply that $\underline{u}=u_{\alpha,\beta,f}$
\deleted{\illegible{1}} normalizes the subalgebra of the matrix algebra
$M_2(\widehat{\mathbb Z})$ \uline{generated} by $\sigma_\infty$
(taking into account the relation
\[
 (99)\qquad \tau_\infty(\sigma_\infty)=-\sigma_\infty
 \qquad\text{in }\mathrm{Gl}(2,\widehat{\mathbb Z}).
\]

If we look at the subgroup of $\mathcal M_{0,3}^{!\sim}$
consisting of the $u$ satisfying the equivalent conditions (93) (94) (98)
\[
 (100)\qquad
 \mathcal M_{0,3}^{!\sim}\left(\tfrac12\right)
 =
 \{u\in\mathcal M_{0,3}^{!\sim}\mid [u,\sigma_\infty]=1\}
 \overset{\sim}{\longrightarrow}\mathcal N_{0,3}^{!\sim}
\]
we find that the composite homomorphism

\scanpage{208}
\mspage{403}

(101)
\diagram{
\begin{tikzcd}[ampersand replacement=\&,column sep=large]
\mathcal N_{0,3}^{!\sim}
 \arrow[r,"\sim"]
 \arrow[rr,bend right=22,"k_{1/2}"']
\&
\mathcal M_{0,3}^{!}(\tfrac12)^\sim
 \arrow[r]
\&
\mathcal M_{0,3}^{!\sim}
 \arrow[r,"{u\longmapsto\underline{u}=(u_{\alpha,\beta,f})}"']
\&
\mathrm{Gl}(2,\widehat{\mathbb Z})'
\\
\unsure{\mathcal N_{0,3}^{!\sim}(\tfrac12)}
\end{tikzcd}
}

\margin{\uline{N.B.} $\mathcal M_{0,3}^{!}(\tfrac12)^\sim$
is stable under $x\longmapsto\sigma_\infty x\sigma_\infty^{-1}$
and under $x\longmapsto\tau_\infty x\tau_\infty^{-1}$
(since $\sigma_\infty$ and $\tau_\infty$
\unsure{commute} with $\sigma_\infty$).}

factors through the said normalizer --- if we prefer to view it in terms of the
group $\mathrm{Gl}(2,\widehat{\mathbb Z})\simeq M_2(\widehat{\mathbb Z})^*$,
rather than of the \uline{algebra} of matrices, it factors through
the \uline{normalizer} of the \uline{maximal torus}
\deleted{(\unsure{formed by})}
\[
 (102)\qquad
 \unsure{T_\alpha|T_\sigma}=
 \{t\sigma_\infty+u\quad
 (t,u\in\unsure{\widehat{\mathbb Z}|\mathbb Z},
 \underbrace{t^2+u^2}_{=\det(t\sigma_\infty+u)}
 \in\widehat{\mathbb Z}^{*}).
\]

I suspect that there is \uline{only the} splitting $k_{1/2}$ of (101)
of the extension $\mathcal M_{0,3}^{!\sim}$
(of $\mathcal N_{0,3}^{!\sim}$ by $\widehat{\Pi}_{0,3}$,
\deleted{\unsure{I think}} or of $\mathcal N_{0,3}^{!\sim}$
by $\widehat{\mathcal G}_{0,3}^{+}$)
that factors through the normalizer.

\margin{Surely bogus:
there are also \unsure{all} the curves \uuline{\unsure{isogenous}}
to the elliptic curve with modular invariant $\unsure{j|J}=0$
(\unsure{corresponding} \deleted{in}
the orbit $\{\tfrac12,2,-1\}$ in $U_{0,3}$ under $\unsure{\mathfrak S_3}$).}

\uline{2° Case}\quad This is the case
\[
 (103)\qquad
 g=1\quad\text{i.e. }\unsure{\alpha_0}=1
 \quad\text{i.e. }\alpha=f
 \quad\text{i.e. }\xi=f\rho(f)^{-1},\ \text{i.e.}
\]
\[
 (104)\qquad [u,\rho]=1.
\]

For an arbitrary $u$, there is a
$\gamma\inserted{\in\widetilde{\widehat{\Pi}}_{0,3}
=\widehat{\Pi}_{0,3}\cdot\{1,\tau_\infty\}}$
\deleted{\unsure{unique}} such that
$u'=\mathrm{int}(\gamma)\circ u=\gamma u$
\unsure{commutes|commuting} thus with $\rho$,
namely $\gamma=\unsure{\alpha_0^{-1}}$,
\deleted{\unsure{We thus have}}\inserted{But} we will pay Attention to the fact that
$\gamma$ is not always in $\widehat{\Pi}_{0,3}$, in fact
\[
 (105)\qquad
 \gamma=\unsure{\alpha_0^{-1}}\in\widehat{\Pi}_{0,3}
 \Longleftrightarrow
 \varepsilon_3(\unsure{u})=1
 \quad\text{i.e. }\mu(u)=1\ (3)
\]

\scanpage{209}

If we denote by
$\unsure{\mathcal M_{0,3}^{!\sim}}(\unsure{-\overline{j}})$
the subgroup of $\mathcal M_{0,3}^{!\sim}$
consisting of the $u$ satisfying (103), (104), its image in
$\unsure{\mathcal N_{0,3}^{!\sim}}$ is therefore
the subgroup of index $2$ kernel of
\[
 u\longmapsto
 \underbrace{\varepsilon_3(\mu(u))}_{\unsure{\sim}}
 =\varepsilon_3(u)
\]
i.e.
\[
 \unsure{(106)}\qquad
 \mathcal N_{0,3}^{!\sim}(-\overline{j})
 =\{u\in\mathcal N_{0,3}^{!\sim}\mid
 \varepsilon_3(u)=1\quad
 \text{i.e. }\mu(u)\equiv1\ (3)\}
\]

\margin{\uline{N.B.} $\mathcal M_{0,3}^{!\sim}(-\overline{j})$
is stable under $x\longmapsto\rho x\rho^{-1}$.}

and we have
\[
 \unsure{(107)}\qquad
 \unsure{\mathcal M_{0,3}^{!\sim}}[-\overline{j}]
 \overset{\sim}{\longrightarrow}
 \mathcal N_{0,3}^{!\sim}[-\overline{j}].
\]

Inverting it we thus find a
\deleted{composite homomorphism}\inserted{splitting}

(108)
\diagram{
\begin{tikzcd}[ampersand replacement=\&,column sep=large]
\mathcal N_{0,3}^{!\sim}(\unsure{-\overline{j}})
 \arrow[r,"\sim"]
 \arrow[rr,bend right=20,"{k(-\overline{j})}"']
\&
\mathcal M_{0,3}^{!\sim}(\unsure{-\overline{j}})
 \arrow[r,hook]
\&
\mathcal M_{0,3}^{!\sim}
\end{tikzcd}
}
,
whose composite with $\mathcal M_{0,3}^{!\sim}$ is a homomorphism which
\deleted{\unsure{factors}}\inserted{factors}
through the maximal torus $T_\rho$ of $\mathrm{Gl}(2,\widehat{\mathbb Z})'$,
formed by the
\deleted{$\unsure{c\rho+d}$}
\[
 \unsure{(109)}\qquad
 \{c\rho+d\mid c,d\in\widehat{\mathbb Z},\
 c^2+cd+d^2\in\widehat{\mathbb Z}^{*}).
\]

I suspect that this is the only splitting of
$\mathcal M_{0,3}^{!\sim}$
(or \unsure{even} of $\mathcal M_{0,3}^{!}$)
which has this property. \unsure{This} follows from the equivalent form
of conditions (103) (104), provided by (79):

\margin{\unsure{bogus}, for \unsure{same} reasons as just now\ldots}

\scanpage{210}
\mspage{404}

\[
 (110)\qquad
 \underline{u}=u_{\alpha,\beta,f}=\mu(c\rho+d)
\]
(N.B. $\delta=\varepsilon_3(\unsure{\mu|u})$ is equal to $1$ in this case).

The $u$ in question correspond to the pairs $(\eta,f)$ satisfying
\[
 (111)\qquad \eta\sigma_\infty(\eta)=1
\]
\[
 \unsure{(112)}\qquad
 (f\rho(f)^{-1})=\eta\ell_1^{\unsure{\nu}}
\]

\margin{which we \unsure{will complete|complete} by $\xi=f\rho(f)^{-1}$.}

\margin{We can also write $\beta$ instead of \unsure{$f$},
$f(=\beta)$ \unsure{determined} modulo
$f\longmapsto\ell_0^{\unsure{n}}f$,
\unsure{id.} for $\alpha,\beta$.}

or again to the pairs $(\alpha,\beta)$ satisfying the loop relation
$(97)=(37)$, which we then complete by $f=\alpha$.

\uline{3° Case}\quad It is the one where
\[
 (113)\qquad
 f\equiv1\ (L_0)
 \quad\text{i.e. }\alpha\equiv\alpha_0\ (L_0\text{ on the left})
 \quad\text{i.e. }\beta\equiv\beta_0\ (L_0\text{ on the left})
\]
or again $u(\varepsilon_0)=\varepsilon_0^\mu$
$(\mu\in\widehat{\mathbb Z}^{*})$ i.e.
\[
 (114)\qquad
 [u,\varepsilon_0]=\ell_0^\nu
 \qquad
 (\nu\in\widehat{\mathbb Z}\ -\ \nu=\tfrac{\mu-1}{2}),
\]
\inserted{i.e. $u$ normalizes
$L_0=\{\ell_0^n\mid n\in\widehat{\mathbb Z}\}$ --- or again
(\unsure{enlarging} $\widehat{\Pi}_{0,3}$)
$u$ normalizes $L'_0=\{\varepsilon_0^n\mid n\in\widehat{\mathbb Z}\}$.}

\margin{\uline{N.B.} $L_0=L'_0\cap\widehat{\Pi}_{0,3}$.}

In this case, we \unsure{remove|can remove} the indeterminacy
$(\alpha,\beta,f)\longmapsto
(\ell_0^n\alpha,\ell_0^n\beta,\ell_0^n f)$
by requiring
\[
 (115)\qquad f=1,
\]
which determines $\alpha,\beta$ uniquely.

\deleted{The set}\inserted{group}
$\mathcal M_{0,3}^{!\sim}[0]$ of the $u\in\mathcal M_{0,3}^{!\sim}$
which normalize $L_0$ is in one-to-one correspondence with the set of pairs
$(\alpha,\beta)$
$(\alpha\in\widetilde{\widehat{\Pi}}_{0,3}
=\widehat{\Pi}_{0,3}\cdot\{1,\tau_\infty\},
\ \beta\in\widehat{\Pi}_{0,3})$ satisfying

\margin{We do not \unsure{lift|lift} Therefore \illegible{2}
by $\mathcal M_{0,3}^{!\sim}(0)$,
\illegible{2} $\mathcal M_{0,3}$ and $\mathcal N_{0,3}$
since it \unsure{does not cover} \illegible{2}
section of the extension \illegible{2} by $\widehat{\Pi}_{0,3}$
--- we \illegible{2} (Proposition \illegible{1})
section $\mathcal M_{0,3}^{!\sim}(0)$ further
\unsure{page 412$^1$}).}

\scanpage{211}

the everlasting loop relation
\[
 \unsure{(115)}\qquad
 (\alpha\rho(\alpha)^{-1})
 =(\beta\sigma_\infty(\beta)^{-1})\ell_1^\nu
 \qquad(\text{for }\nu\in\widehat{\mathbb Z}\text{ suitable})
\]
\inserted{(by $(\alpha,\beta)\longmapsto
\unsure{g|\xi}=\alpha\rho(\alpha)^{-1};
\unsure{h|\eta}=\beta\sigma_\infty(\beta)^{-1}$)}
or again to the pairs $(g,h)$
\inserted{(\unsure{or} $(\xi,\eta)$)}
of elements of $\widehat{\Pi}_{0,3}$ satisfying
\[
 \begin{gathered}
 (116)\quad (=(44)+(46))\\
 g\rho(g)\rho^2(g)=1,\qquad
 \unsure{h}\sigma_\infty(\unsure{h})=1,\qquad
 g=\unsure{h}\ell_1^\nu\quad
 (\nu\in\widehat{\mathbb Z}\text{ suitable}),
 \end{gathered}
\]

\margin{In all this, I drop the ``preliminary
condition'' on the generating relations\ldots}

to the pair \deleted{$(\alpha,\beta)$ \unsure{as}}
$(g,h)$ corresponds the automorphism $u$ defined by
\[
 (117)=(23)\qquad
 u(\sigma_\infty)=\unsure{h}\sigma_\infty,\qquad
 u(\rho)=g\rho
\]
or again by
\[
 (118)\qquad
 \unsure{u}(\sigma_\infty)=\mathrm{int}(\alpha)\sigma_\infty,\qquad
 u(\rho)=\mathrm{int}(\beta)\rho.
\]

In the group $\mathcal M_{0,3}^{!\sim}[0]$, there is the subgroup
\[
 (119)\qquad
 \mathcal M_{0,3}^{!\sim}[0]\cap\widehat{\Pi}_{0,3}\simeq L_0
\]
the inclusion
\[
 (120)\qquad
 \widehat{\mathbb Z}\overset{\ell_0}{\underset{\sim}{\longrightarrow}}
 L_0\lhook\joinrel\longrightarrow\mathcal M_{0,3}^{!\sim}
\]
being described by
\[
 (121)\qquad
 \lambda=\ell_0^n\longmapsto
 (\alpha=\unsure{\lambda|\ell_0^n},
 \beta=\unsure{\lambda|\ell_0^n})
\]
\deleted{\unsure{i.e.}\ \illegible{2}\ $(\alpha=\lambda,\beta=$}

or again by
\[
 (122)\qquad
 \lambda=\ell_0^{\unsure{n}}\longmapsto
 (f=\ell_0^n\ell_1^{-n},g=\ell_0^n\ell_1^{-n})
\]
\inserted{i.e. $f=g=\lambda\sigma_\infty(\unsure{\lambda^{-1}})
=\lambda\rho(\unsure{\lambda})^{-1}$,}
multiplication \deleted{by} on the left by \unsure{the}
$\lambda=\ell_0^n$ in $\mathcal M_{0,3}^{!\sim}[0]$
being \unsure{made explicit|made explicit} by
\[
 (123)\qquad
 (\alpha,\beta)\longmapsto(\ell_0^n\alpha,\ell_0^n\beta)
\]
\[
 (124)\qquad
 (\xi,\eta)\longmapsto
 \ell_0^n\xi\ell_1^{-n},\ \ell_0^n\eta\ell_1^{-n}).
\]

\scanpage{212}
\mspage{405}

Since we have an isomorphism
\[
\mathcal M_{0,3}^{!\sim}[0]/L_0
\xrightarrow{\sim}\mathcal N_{0,3}^{!\sim},
\]
we find the

\uline{Corollary} The group $\mathcal N_{0,3}^{!\sim}$ is in canonical bijection
with the set of pairs
$(\alpha,\beta)\in\Pi_{0,3}^{\sim\wedge}\times\Pi_{0,3}^{\wedge}$
satisfying the loop relation (97)
\[
\alpha\rho(\alpha)^{-1}
=\bigl(\beta\sigma_\infty(\beta)^{-1}\bigr)\ell_1^\nu
\qquad(\text{for }\nu\in\widehat{\mathbb Z}\text{ suitable}),
\]
\uline{modulo} the transformations
$(\alpha,\beta)\mapsto(\ell_0^{\unsure{n|m}}\alpha,\ell_0^{\unsure{n|m}}\beta)$.

The group law \inserted{(of $\mathcal M_{0,3}^{!\sim}(0)$ or of its quotient
$\mathcal N_{0,3}^{!\sim}$)} is made explicit in terms of the $(g,h)$, or of the
$(\alpha,\beta)$, by the ``cocycle formulas'' (52): if $u',u$
\inserted{and $u''=u'\cdot u$} are given by
$\unsure{(g',h')},(g,h),(g'',h'')$, or by
$(\alpha',\beta'),(\alpha,\beta)$ and $(\alpha'',\beta'')$, we will have
\[
\tag{125}
\left\{
\begin{aligned}
g''&=u'(g)g'&h''&=u'(h)h'\\
\alpha''&=u'(\alpha)\alpha'&\beta''&=u'(\beta)\beta'.
\end{aligned}
\right.
\]

\uline{Remark on the notation} $\mathcal M_{0,3}^{!\sim}(q)$,
for $q=\frac12,-\overline j,0$.
In fact for every \uline{algebraic} point
\[
\tag{126}
q\in\mathcal U_{0,3}(\mathbb C)
=\mathbb C\setminus\{0,1\}
=\mathbb P^1_{\unsure{\mathbb C|\mathbb Q}}(\mathbb C)
\setminus\{0,1,\infty\}
\]
we expect \deleted{\unsure{to find}} a partial splitting
\[
\tag{127}
k_q:\unsure{\mathcal N|\mathcal M}_{0,3}^{!\sim}(q)\longrightarrow\mathcal M_{0,3}^{!\sim},
\qquad\operatorname{Im}k_q\overset{\text{definition}}=\mathcal M_{0,3}^{!\sim}(q)
\]
of the extension $\mathcal M_{0,3}^{!\sim}$ of
$\mathcal N_{0,3}^{!\sim}$ by $\Pi_{0,3}^{\wedge}$, or

\margin{In fact, \unsure{in a way} \illegible{1} more
\unsure{precise|precise}, $k_q$ will only be \unsure{determined}
up to an inner \unsure{automorphism}...}